\begin{textsample}{POS Dim 1 – 2015 – Score 20.00 – 2015/t001601.txt}  \label{ex:f1_pos_041}
I ' d like to take you on \textbf{the} epic quest of \textbf{the} Rosetta spacecraft.
To escort and land \textbf{the} probe on a comet, this has been my passion for \textbf{the} past two years.
In order to do that, I need to explain to you something about \textbf{the} origin of \textbf{the} solar system.
When we go back four and a half billion years, there was a cloud of gas and dust.
In \textbf{the} center of this cloud, our sun formed and ignited.
Along with that, what we now know as \textbf{planets}, comets and \textbf{asteroids} formed.
What then happened, according to theory, is that when \textbf{the} Earth had cooled down a bit after its formation, comets massively impacted \textbf{the} Earth and delivered water to Earth.
They probably also delivered complex \textbf{organic} material to Earth, and that may have bootstrapped \textbf{the} emergence of life.
You can compare this to having to solve a 250-piece puzzle and not a 2, 000-piece puzzle.
Afterwards, \textbf{the} big \textbf{planets} like Jupiter and Saturn, they were not in their place where they are now, and they interacted gravitationally, and they swept \textbf{the} whole interior of \textbf{the} solar system clean, and what we now know as comets ended up in something called \textbf{the} Kuiper Belt, which is a belt of \textbf{objects} beyond \textbf{the} orbit of Neptune.
And sometimes these \textbf{objects} run into each other, and they gravitationally deflect, and then \textbf{the} gravity of Jupiter pulls them back into \textbf{the} solar system.
And they then become \textbf{the} comets as we see them in \textbf{the} sky. \textbf{The} important thing here to note is that in \textbf{the} meantime, \textbf{the} four and a half billion years, these comets have been sitting on \textbf{the} outside of \textbf{the} solar system, and haven ' t changed — deep, frozen versions of our solar system.
In \textbf{the} sky, they look like this.
We know them for their tails.
There are actually two tails.
One is a dust tail, which is blown away by \textbf{the} solar wind. \textbf{The} other one is an ion tail, which is charged particles, and they follow \textbf{the} magnetic field in \textbf{the} solar system.
There ' s \textbf{the} coma, and then there is \textbf{the} nucleus, which here is too small to see, and you have to remember that in \textbf{the} case of Rosetta, \textbf{the} spacecraft is in that center pixel.
We are only 20, 30, 40 kilometers away from \textbf{the} comet.
So what ' s important to remember?
Comets contain \textbf{the} original material from which our solar system was formed, so they ' re ideal to study \textbf{the} components that were present at \textbf{the} time when Earth, and life, started.
Comets are also suspected of having brought \textbf{the} elements which may have bootstrapped life.
In 1983, ESA set up its long-term Horizon 2000 program, which contained one cornerstone, which would be a mission to a comet.
In parallel, a small mission to a comet, what you see here, Giotto, was launched, and in 1986, flew by \textbf{the} comet of Halley with an armada of other spacecraft.
From \textbf{the} results of that mission, it became immediately clear that comets were ideal bodies to study to understand our solar system.
And thus, \textbf{the} Rosetta mission was approved in 1993, and originally it was supposed to be launched in 2003, but a problem arose with an Ariane rocket.
However, our P.
R. department, in its enthusiasm, had already made 1, 000 Delft \textbf{Blue} \textbf{plates} with \textbf{the} name of \textbf{the} wrong comets.
So I ' ve never had to buy any china since.
That ' s \textbf{the} positive part.
Once \textbf{the} whole problem was solved, we left Earth in 2004 to \textbf{the} newly selected comet, Churyumov-Gerasimenko.
This comet had to be specially selected because A, you have to be able to get to it, and B, it shouldn ' t have been in \textbf{the} solar system too long.
This particular comet has been in \textbf{the} solar system since 1959.
That ' s \textbf{the} first time when it was deflected by Jupiter, and it got close enough to \textbf{the} sun to start changing.
So it ' s a very fresh comet.
Rosetta made a few historic firsts.
It ' s \textbf{the} first satellite to orbit a comet, and to escort it throughout its whole tour through \textbf{the} solar system — closest approach to \textbf{the} sun, as we will see in August, and then away again to \textbf{the} exterior.
It ' s \textbf{the} first ever landing on a comet.
We actually orbit \textbf{the} comet using something which is not normally done with spacecraft.
Normally, you look at \textbf{the} sky and you know where you point and where you are.
In this case, that ' s not enough.
We navigated by looking at landmarks on \textbf{the} comet.
We recognized features — boulders, craters — and that ' s how we know where we are respective to \textbf{the} comet.
And, of course, it ' s \textbf{the} first satellite to go beyond \textbf{the} orbit of Jupiter on solar cells.
Now, this sounds more heroic than it actually is, because \textbf{the} technology to use radio isotope thermal generators wasn ' t available in Europe at that time, so there was no choice.
But these solar \textbf{arrays} are big.
This is one \textbf{wing}, and these are not specially selected small people.
They ' re just like you and me.
We have two of these \textbf{wings}, 65 square meters.
Now later on, of course, when we got to \textbf{the} comet, you find out that 65 square meters of sail close to a body which is outgassing is not always a very handy choice.
Now, how did we get to \textbf{the} comet?
Because we had to go there for \textbf{the} Rosetta scientific objectives very far away — four times \textbf{the} distance of \textbf{the} Earth to \textbf{the} sun — and also at a much higher velocity than we could achieve with fuel, because we ' d have to take six times as much fuel as \textbf{the} whole spacecraft weighed.
So what do you do?
You use gravitational flybys, slingshots, where you pass by a \textbf{planet} at very low \textbf{altitude}, a few thousand kilometers, and then you get \textbf{the} velocity of that \textbf{planet} around \textbf{the} sun for free.
We did that a few times.
We did Earth, we did Mars, we did twice Earth again, and we also flew by two \textbf{asteroids}, Lutetia and Steins.
Then in 2011, we got so far from \textbf{the} sun that if \textbf{the} spacecraft got into trouble, we couldn ' t actually save \textbf{the} spacecraft anymore, so we went into hibernation.
Everything was switched off except for one clock.
Here you see in white \textbf{the} trajectory, and \textbf{the} way this works.
You see that from \textbf{the} circle where we started, \textbf{the} white line, actually you get more and more and more elliptical, and then finally we approached \textbf{the} comet in May 2014, and we had to start doing \textbf{the} rendezvous maneuvers.
On \textbf{the} way there, we flew by Earth and we took a few pictures to test our cameras.
This is \textbf{the} moon rising over Earth, and this is what we now call a selfie, which at that time, by \textbf{the} way, that word didn ' t exist.
It ' s at Mars.
It was taken by \textbf{the} CIVA camera.
That ' s one of \textbf{the} cameras on \textbf{the} lander, and it just looks under \textbf{the} solar \textbf{arrays}, and you see \textbf{the} \textbf{planet} Mars and \textbf{the} solar \textbf{array} in \textbf{the} distance.
Now, when we got out of hibernation in January 2014, we started arriving at a distance of two million kilometers from \textbf{the} comet in May.
However, \textbf{the} velocity \textbf{the} spacecraft had was much too fast.
We were going 2, 800 kilometers an hour faster than \textbf{the} comet, so we had to \textbf{brake}.
We had to do eight maneuvers, and you see here, some of them were really big.
We had to \textbf{brake} \textbf{the} first one by a few hundred kilometers per hour, and actually, \textbf{the} duration of that was seven hours, and it used 218 kilos of fuel, and those were seven nerve-wracking hours, because in 2007, there was a leak in \textbf{the} system of \textbf{the} propulsion of Rosetta, and we had to close off a branch, so \textbf{the} system was actually operating at a pressure which it was never designed or qualified for.
Then we got in \textbf{the} vicinity of \textbf{the} comet, and these were \textbf{the} first pictures we saw. \textbf{The} true comet rotation period is 12 and a half hours, so this is accelerated, but you will understand that our flight dynamics engineers thought, this is not going to be an easy thing to land on.
We had hoped for some kind of spud-like thing where you could easily land.
But we had one hope : maybe it was smooth.
No.
That didn ' t work either.
So at that point in time, it was clearly unavoidable : we had to map this body in all \textbf{the} detail you could get, because we had to find an area which is 500 meters in \textbf{diameter} and flat.
Why 500 meters?
That ' s \textbf{the} error we have on landing \textbf{the} probe.
So we went through this process, and we mapped \textbf{the} comet.
We used a \textbf{technique} called photoclinometry.
You use shadows thrown by \textbf{the} sun.
What you see here is a \textbf{rock} sitting on \textbf{the} \textbf{surface} of \textbf{the} comet, and \textbf{the} sun shines from above.
From \textbf{the} shadow, we, with our brain, can immediately determine roughly what \textbf{the} shape of that \textbf{rock} is.
You can program that in a \textbf{computer}, you then cover \textbf{the} whole comet, and you can map \textbf{the} comet.
For that, we flew special trajectories starting in August.
First, a \textbf{triangle} of 100 kilometers on a side at 100 kilometers ' distance, and we repeated \textbf{the} whole thing at 50 kilometers.
At that time, we had seen \textbf{the} comet at all kinds of angles, and we could use this \textbf{technique} to map \textbf{the} whole thing.
Now, this led to a \textbf{selection} of landing sites.
This whole process we had to do, to go from \textbf{the} mapping of \textbf{the} comet to actually finding \textbf{the} final landing site, was 60 days.
We didn ' t have more.
To give you an idea, \textbf{the} average Mars mission takes hundreds of scientists for years to meet about where shall we go?
We had 60 days, and that was it.
We finally selected \textbf{the} final landing site and \textbf{the} commands were prepared for Rosetta to launch Philae. \textbf{The} way this works is that Rosetta has to be at \textbf{the} right point in space, and aiming towards \textbf{the} comet, because \textbf{the} lander is passive. \textbf{The} lander is then pushed out and moves towards \textbf{the} comet.
Rosetta had to turn around to get its cameras to actually look at Philae while it was departing and to be able to communicate with it.
Now, \textbf{the} landing duration of \textbf{the} whole trajectory was seven hours.
Now do a simple calculation : if \textbf{the} velocity of Rosetta is off by one centimeter per second, seven hours is 25, 000 seconds.
That means 252 meters wrong on \textbf{the} comet.
So we had to know \textbf{the} velocity of Rosetta much better than one centimeter per second, and its location in space better than 100 meters at 500 million kilometers from Earth.
That ' s no mean feat.
Let me quickly take you through some of \textbf{the} science and \textbf{the} instruments.
I won ' t bore you with all \textbf{the} details of all \textbf{the} instruments, but it ' s got everything.
We can sniff gas, we can measure dust particles, \textbf{the} shape of them, \textbf{the} composition, there are magnetometers, everything.
This is one of \textbf{the} results from an instrument which measures gas density at \textbf{the} position of Rosetta, so it ' s gas which has left \textbf{the} comet. \textbf{The} bottom graph is September of last year.
There is a long-term variation, which in itself is not surprising, but you see \textbf{the} sharp peaks.
This is a comet day.
You can see \textbf{the} effect of \textbf{the} sun on \textbf{the} evaporation of gas and \textbf{the} fact that \textbf{the} comet is rotating.
So there is one spot, apparently, where there is a lot of stuff coming from, it gets heated in \textbf{the} Sun, and then cools down on \textbf{the} back side.
And we can see \textbf{the} density variations of this.
These are \textbf{the} gases and \textbf{the} \textbf{organic} compounds that we already have measured.
You will see it ' s an impressive list, and there is much, much, much more to come, because there are more measurements.
Actually, there is a conference going on in Houston at \textbf{the} moment where many of these results are presented.
Also, we measured dust particles.
Now, for you, this will not look very impressive, but \textbf{the} scientists were thrilled when they saw this.
Two dust particles : \textbf{the} right one they call Boris, and they shot it with tantalum in order to be able to analyze it.
Now, we found sodium and magnesium.
What this tells you is this is \textbf{the} concentration of these two materials at \textbf{the} time \textbf{the} solar system was formed, so we learned things about which materials were there when \textbf{the} \textbf{planet} was made.
Of course, one of \textbf{the} important elements is \textbf{the} imaging.
This is one of \textbf{the} cameras of Rosetta, \textbf{the} OSIRIS camera, and this actually was \textbf{the} cover of Science magazine on January 23 of this year.
Nobody had expected this body to look like this.
Boulders, \textbf{rocks} — if anything, it looks more like \textbf{the} Half Dome in Yosemite than anything else.
We also saw things like this : dunes, and what look to be, on \textbf{the} righthand side, wind-blown shadows.
Now we know these from Mars, but this comet doesn ' t have an atmosphere, so it ' s a bit difficult to create a wind-blown shadow.
It may be local outgassing, stuff which goes up and comes back.
We don ' t know, so there is a lot to investigate.
Here, you see \textbf{the} same image twice.
On \textbf{the} left-hand side, you see in \textbf{the} middle a pit.
On \textbf{the} \textbf{right-hand} side, if you carefully look, there are three jets coming out of \textbf{the} bottom of that pit.
So this is \textbf{the} activity of \textbf{the} comet.
Apparently, at \textbf{the} bottom of these pits is where \textbf{the} active regions are, and where \textbf{the} material evaporates into space.
There is a very intriguing crack in \textbf{the} neck of \textbf{the} comet.
You see it on \textbf{the} \textbf{right-hand} side.
It ' s a kilometer long, and it ' s two and a half meters wide.
Some people suggest that actually, when we get close to \textbf{the} sun, \textbf{the} comet may split in two, and then we ' ll have to choose, which comet do we go for? \textbf{The} lander — again, lots of instruments, mostly comparable except for \textbf{the} things which hammer in \textbf{the} ground and drill, etc.
But much \textbf{the} same as Rosetta, and that is because you want to compare what you find in space with what you find on \textbf{the} comet.
These are called ground truth measurements.
These are \textbf{the} landing descent images that were taken by \textbf{the} OSIRIS camera.
You see \textbf{the} lander getting further and further away from Rosetta.
On \textbf{the} top right, you see an image taken at 60 meters by \textbf{the} lander, 60 meters above \textbf{the} \textbf{surface} of \textbf{the} comet. \textbf{The} boulder there is some 10 meters.
So this is one of \textbf{the} last images we took before we landed on \textbf{the} comet.
Here, you see \textbf{the} whole sequence again, but from a different perspective, and you see three blown-ups from \textbf{the} bottom-left to \textbf{the} middle of \textbf{the} lander traveling over \textbf{the} \textbf{surface} of \textbf{the} comet.
Then, at \textbf{the} top, there is a before and an after image of \textbf{the} landing. \textbf{The} only problem with \textbf{the} after image is, there is no lander.
But if you carefully look at \textbf{the} \textbf{right-hand} side of this image, we saw \textbf{the} lander still there, but it had bounced.
It had departed again.
Now, on a bit of a comical note here is that originally Rosetta was designed to have a lander which would bounce.
That was \textbf{discarded} because it was way too expensive.
Now, we forgot, but \textbf{the} lander knew.
During \textbf{the} first bounce, in \textbf{the} magnetometers, you see here \textbf{the} data from them, from \textbf{the} three axes, x, y and z.
Halfway through, you see a red line.
At that red line, there is a change.
What happened, apparently, is during \textbf{the} first bounce, somewhere, we hit \textbf{the} edge of a crater with one of \textbf{the} legs of \textbf{the} lander, and \textbf{the} rotation velocity of \textbf{the} lander changed.
So we ' ve been rather lucky that we are where we are.
This is one of \textbf{the} iconic images of Rosetta.
It ' s a man-made \textbf{object}, a leg of \textbf{the} lander, standing on a comet.
This, for me, is one of \textbf{the} very best images of space science I have ever seen.
One of \textbf{the} things we still have to do is to actually find \textbf{the} lander. \textbf{The} \textbf{blue} area here is where we know it must be.
We haven ' t been able to find it yet, but \textbf{the} \textbf{search} is continuing, as are our efforts to start getting \textbf{the} lander to work again.
We listen every day, and we hope that between now and somewhere in April, \textbf{the} lander will wake up again. \textbf{The} findings of what we found on \textbf{the} comet : This thing would float in water.
It ' s half \textbf{the} density of water.
So it looks like a very big \textbf{rock}, but it ' s not. \textbf{The} activity increase we saw in June, July, August last year was a four-fold activity increase.
By \textbf{the} time we will be at \textbf{the} sun, there will be 100 kilos a second leaving this comet : gas, dust, whatever.
That ' s 100 million kilos a day.
Then, finally, \textbf{the} landing day.
I will never forget — absolute madness, 250 TV crews in Germany. \textbf{The} BBC was interviewing me, and another TV crew who was following me all day were filming me being interviewed, and it went on like that for \textbf{the} whole day. \textbf{The} Discovery Channel crew actually caught me when leaving \textbf{the} control room, and they asked \textbf{the} right question, and man, I got into tears, and I still feel this.
For a month and a half, I couldn ' t think about landing day without crying, and I still have \textbf{the} emotion in me.
With this image of \textbf{the} comet, I would like to leave you.
Thank you.

% matched lemmas: altitude, array, asteroid, blue, brake, computer, diameter, discard, object, organic, planet, plate, right-hand, rock, search, selection, surface, technique, the, triangle, wing
\end{textsample}
